% FONTE TEMA https://github.com/matze/mtheme
\documentclass[aspectratio=1610]{beamer}
%\documentclass[aspectratio=1610, handout]{beamer}
\usepackage[utf8]{inputenc}
\usepackage{ragged2e}
\usepackage{xcolor}
\usepackage[italian]{babel}
\usepackage{multirow}
\usepackage{silence}
\usepackage{listings}
\definecolor{verde}{rgb}{0,0.6,0}
\definecolor{grigio}{rgb}{0.5,0.5,0.5}
\definecolor{rosa}{rgb}{0.58,0,0.82}
\definecolor{backgroud}{rgb}{0.95,0.95,0.92}
\lstdefinestyle{stilepython}{
    backgroundcolor=\color{backgroud},
    commentstyle=\color{verde},
    keywordstyle=\color{magenta},
    numberstyle=\tiny\color{grigio},
    stringstyle=\color{rosa},
    basicstyle=\ttfamily\footnotesize,
    breakatwhitespace=false,
    numbers=left,
    numbersep=5pt,
    showspaces=false,
    showstringspaces=false,
    showtabs=false,
    literate=
    {à}{{\'a}}1
    {è}{{\'e}}1
    {ì}{{\'i}}1
    {ò}{{\'o}}1
    {ù}{{\'u}}1
}
\lstset{style=stilepython}

\WarningFilter{beamer}{}
\WarningFilter{metropolis}{}
\usetheme[progressbar=frametitle,titleformat=smallcaps]{metropolis}
\setbeamertemplate{frame numbering}[fraction]
\setbeamercovered{dynamic}
\definecolor{rosso}{RGB}{255, 0, 0}
\definecolor{giallo}{RGB}{254,212,23}
\hypersetup{colorlinks=true,linkcolor=black,urlcolor=rosso}
\setbeamercolor{palette primary}{fg=black, bg=giallo}
\setbeamercolor{background canvas}{bg=white}
\setbeamercolor{normal text}{fg=black}
\setbeamercolor{progress bar}{fg=rosso}
\setbeamercolor{framesubtitle}{fg=rosso}
\setbeamercolor{normal text .dimmed}{fg=giallo}
\setbeamercolor{block title alerted}{fg=rosso, bg=giallo}
\setbeamerfont{caption}{size=\tiny}
\setbeamerfont{caption name}{size=\tiny}
\setlength{\abovecaptionskip}{0pt}
\makeatletter
\metroset{block=fill}
\setlength{\metropolis@progressinheadfoot@linewidth}{1pt} 
\setlength{\metropolis@progressonsectionpage@linewidth}{1pt}
\setlength{\metropolis@titleseparator@linewidth}{1pt}
\makeatother

\title{STRINGHE E LISTE}
\subtitle{Metodi principali e strutture dati di base}
\date{}
\institute{\textit{
        Fonti:
        \begin{itemize}
            \item[-] \href{https://docs.python.org/it/3.15/builtins/stdtypes.html\#string-methods}{Metodi delle Stringhe - Documentazione Ufficiale Python}
            \item[-] \href{https://docs.python.org/it/3.15/tutorial/introduction.html\#lists}{Definizione di Lista - Documentazione Ufficiale Python}
            \item[-] \href{https://docs.python.org/it/3/tutorial/datastructures.html}{Strutture Dati - Documentazione Ufficiale Python} 
        \end{itemize}
    }
}

\begin{document}

\begin{frame}[plain, noframenumbering]
    \titlepage
\end{frame}

\section{STRINGHE}

\begin{frame}[fragile]{METODI PRINCIPALI DELLE STRINGHE}
    \begin{alertblock}{DEFINIZIONE}
        \begin{minipage}{0.98\linewidth}
            \justifying
            Le stringhe in Python sono sequenze di caratteri immutabili. Esse supportano diversi metodi 
            che permettono di manipolare e analizzare il contenuto delle stringhe stesse. I metodi per essere utilizzati 
            devono essere invocati su un oggetto stringa, utilizzando la sintassi:
\begin{lstlisting}[language=Python]
 stringa.metodo(parametro1, parametro2, ...)
\end{lstlisting}
\pause
Esempio:
\begin{lstlisting}[language=Python]
 testo = "Space Invaders"
 lista = testo.split()
 print(lista)  # Output: ['Space', 'Invaders']
\end{lstlisting}
        \end{minipage}
    \end{alertblock}
\end{frame}

\begin{frame}{METODI PRINCIPALI DELLE STRINGHE}
    \begin{alertblock}{METODI PRINCIPALI}
        \begin{minipage}{0.98\linewidth}
            \justifying
            \begin{itemize}
                \item \textbf{split("x")}: Divide la stringa in una lista di sotto-stringhe, utilizzando "x" come delimitatore.
                \pause
                \item \textbf{count("x")}: Conta quante volte il carattere o la sotto-stringa "x" appare nella stringa. 
                \pause
                \item \textbf{lower()}: Restituisce una nuova stringa con tutti i caratteri convertiti in minuscolo.
                \pause
                \item \textbf{upper()}: Restituisce una nuova stringa con tutti i caratteri convertiti in maiuscolo.
                \pause
                \item \textbf{strip()}: Rimuove gli spazi bianchi iniziali e finali dalla stringa.
            \end{itemize}
        \end{minipage}
    \end{alertblock}
\end{frame}

\begin{frame}{METODI PRINCIPALI DELLE STRINGHE}
    \begin{alertblock}{METODI PRINCIPALI}
        \begin{minipage}{0.98\linewidth}
            \justifying
            \begin{itemize}
                \item \textbf{endswith("x")}: Restituisce True se la stringa termina con "x", altrimenti False.
                \pause
                \item \textbf{startswith("x")}: Restituisce True se la stringa inizia con "x", altrimenti False.
                \pause
                \item \textbf{isdecimal()}: Restituisce True se tutti i caratteri della stringa sono numeri decimali, altrimenti False.
                \pause
                \item \textbf{islower()}: Restituisce True se tutti i caratteri della stringa sono minuscoli, altrimenti False.
                \pause
                \item \textbf{isupper()}: Restituisce True se tutti i caratteri della stringa sono maiuscoli, altrimenti False.
            \end{itemize}
        \end{minipage}
    \end{alertblock}
\end{frame}

\begin{frame}{FUNZIONI BUILT-IN PER LE STRINGHE}
    \begin{alertblock}{FUNZIONI PRINCIPALI}
        \begin{minipage}{0.98\linewidth}
            \justifying
            \begin{itemize}
                \item \textbf{str(x)}: Converte l'oggetto x in una stringa.
                \pause
                \item \textbf{len("x")}: Restituisce la lunghezza della stringa x.
                \pause
                \item \textbf{ord("c")}: Restituisce il valore numerico del carattere c.
                \pause
                \item \textbf{chr(x)}: Restituisce il carattere corrispondente al valore numerico x.
            \end{itemize}
        \end{minipage}
    \end{alertblock}
\end{frame}

\begin{frame}[fragile]{STRINGHE E LISTE}
    \begin{alertblock}{DEFINIZIONE}
        \begin{minipage}{0.98\linewidth}
            \justifying
            Le stringhe possono essere viste come liste di caratteri; per poterle manipolare, 
            è possibile accedere ai singoli caratteri utilizzando gli indici.\\
            \textbf{Un indice è un numero intero all'interno di parentesi quadre che rappresenta la posizione di un 
            elemento all'interno di una sequenza}. Gli indici \textbf{partono da 0}.\\
\pause
Esempio:
\begin{lstlisting}[language=Python]
 testo = "cyberpunk"
 carattere = testo[0]  # carattere conterrà 'c'
 carattere = testo[1]  # carattere conterrà 'y'
\end{lstlisting}
        \end{minipage}
    \end{alertblock}
\end{frame}

\begin{frame}[fragile]{STRINGHE E LISTE}
    \begin{alertblock}{DEFINIZIONE}
        \begin{minipage}{0.98\linewidth}
            \justifying
            Lo \textbf{Slicing} permette di estrarre una porzione di una stringa utilizzando la sintassi:\\
            \center{\textbf{stringa[ inizio : fine : passo ]}}
            \pause
            \begin{itemize}
                \item \textbf{inizio}: indice del primo carattere da includere nello slice (inclusivo).
                \pause
                \item \textbf{fine}: indice del primo carattere da escludere dallo slice (esclusivo).
                \pause
                \item \textbf{passo}: intervallo tra gli indici da includere nello slice.
            \end{itemize}
\pause
\raggedright
Esempio:
\begin{lstlisting}[language=Python]
 testo = "cyberpunk"
 substr1 = testo[0:5]  # substr1 conterrà 'cyber'
 substr2 = testo[0:10:2]  # substr2 conterrà 'cbruk'
\end{lstlisting}
        \end{minipage}
    \end{alertblock}
\end{frame}

\section{LISTE}

\begin{frame}[fragile]{LISTE}
    \begin{alertblock}{DEFINIZIONE}
        \begin{minipage}{0.98\linewidth}
            \justifying
            Una lista è una struttura dati che può contenere una sequenza di elementi, che possono essere 
            di qualsiasi tipo (numeri, stringhe, ecc.). Le \textbf{liste sono mutabili}, il che 
            significa che è possibile modificare i loro elementi dopo la creazione.\\
            Una lista può essere creata utilizzando le parentesi quadre \textbf{[ ]} e separando 
            gli elementi con le virgole oppure utilizzando la funzione \textbf{list()}.\\
            \pause
            Esempio:
\begin{lstlisting}[language=Python]
 lista1 = [1 , 2 , 3 , 4 , 5]
 lista2 = ["a" , "b" , "c"]
 lista3 = ["Age of Empires" , 1997 , True]
 lista4 = list(range(5))  # lista4 conterrà [0, 1, 2, 3, 4]
 lista5 = list("GTA")  # lista5 conterrà ['G', 'T', 'A']
\end{lstlisting}
        \end{minipage}
    \end{alertblock}
\end{frame}

\begin{frame}[fragile]{LISTE}
    \begin{alertblock}{ESEMPIO}
        \begin{minipage}{0.98\linewidth}
\begin{lstlisting}[language=Python]
 lista = list(map(int,input("Inserisci numeri interi: ").split()))
\end{lstlisting}
\pause
\begin{enumerate}
    \item La funzione \textbf{input()} legge una riga di input dall'utente.
    \pause
    \item Il metodo \textbf{split()} divide la stringa di input in una lista di sotto-stringhe, utilizzando lo spazio come delimitatore.
    \pause
    \item La funzione \textbf{map()} applica la funzione \textbf{int()} a ciascun elemento della lista di sotto-stringhe, convertendoli in numeri interi.
    \pause
    \item La funzione \textbf{list()} converte l'oggetto mappa risultante in una lista di numeri interi.
    \pause
    \item la lista risultante viene assegnata alla variabile \textbf{lista}.
\end{enumerate}
        \end{minipage}
    \end{alertblock}
\end{frame}

\begin{frame}{FUNZIONI BUILT-IN PER LE LISTE}
    \begin{alertblock}{FUNZIONI PRINCIPALI}
        \begin{minipage}{0.98\linewidth}
            \justifying
            \begin{itemize}
                \item \textbf{len(lista)}: Restituisce il numero di elementi nella lista.
                \pause
                \item \textbf{min(lista)}: Restituisce il valore minimo nella lista.
                \pause
                \item \textbf{max(lista)}: Restituisce il valore massimo nella lista.
                \pause
                \item \textbf{sum(lista)}: Restituisce la somma di tutti gli elementi numerici nella lista.
            \end{itemize}
        \end{minipage}
    \end{alertblock}
\end{frame}

\begin{frame}{METODI PRINCIPALI PER LE LISTE}
    \begin{alertblock}{METODI PRINCIPALI}
        \begin{minipage}{0.98\linewidth}
            \justifying
            \begin{itemize}
                \item \textbf{append(x)}: Aggiunge l'elemento x alla fine della lista.
                \pause
                \item \textbf{insert(indice, x)}: Inserisce l'elemento x nella posizione indice della lista.
                \pause
                \item \textbf{remove(x)}: Rimuove la prima occorrenza dell'elemento x dalla lista.
                \pause
                \item \textbf{pop(i)}: Rimuove e restituisce l'elemento nella posizione i della lista. Se i non è specificato, rimuove e restituisce l'ultimo elemento.
                \pause
                \item \textbf{clear()}: Rimuove tutti gli elementi dalla lista.
                \pause
                \item \textbf{sort()}: Ordina gli elementi della lista in ordine crescente.
                \pause
                \item \textbf{reverse()}: Inverte l'ordine degli elementi nella lista.
            \end{itemize}
        \end{minipage}
    \end{alertblock}
\end{frame}

\begin{frame}[fragile]{LISTE}
    \begin{alertblock}{ESEMPIO}
        \begin{minipage}{0.98\linewidth}
            Un'app di running registra i tempi in minuti delle sessioni di corsa di un atleta. 
            Lo script deve inserire una nuova sessione e calcolare il tempo corso totale.
            Infine, invertire l'ordine della lista per mostrare la cronologia dalla sessione 
            più recente alla più vecchia.
\begin{lstlisting}[language=Python]
tempi = [45, 30, 50, 30, 60, 25]
tempi.append(35)
totale = sum(tempi)
tempi.reverse()
print("---Tempo totale allenamenti---")
print(f"{totale} minuti in {len(tempi)} sessioni")
print("Cronologia allenamenti (dal più recente):", tempi)
\end{lstlisting}
        \end{minipage}
    \end{alertblock}
\end{frame}

\section{DIZIONARI}

\begin{frame}[fragile]{DIZIONARI}
    \begin{alertblock}{DEFINIZIONE}
        \begin{minipage}{0.98\linewidth}
            \justifying
            Un dizionario è una struttura dati che memorizza \textbf{coppie chiave-valore}. 
            \textbf{Ogni chiave deve essere unica e immutabile, mentre i valori possono essere 
            di qualsiasi tipo}. I dizionari sono mutabili, il che significa che è possibile 
            aggiungere, modificare o rimuovere coppie chiave-valore dopo la creazione del dizionario.\\
            La sintassi per creare un dizionario è la seguente:
            \center{\textbf{dizionario = \{chiave1: valore1 , chiave2: valore2 , ...\}}}\\
            \pause
            \raggedright
            Esempi:
\begin{lstlisting}[language=Python]
 tel = {"Ada": 3481234567, "Bea": 3479876543, "Eva": 3456789123}
 tel = dict(ada=3481234567, Bea=3479876543, Eva=3456789123)
 tel["Ada"] # Restituisce 3481234567
 tel["Mia"] = 3498765432 # Aggiunge una nuova coppia chiave-valore
 list(tel) # Restituisce la lista delle chiavi del dizionario
 sorted(tel) # Restituisce la lista delle chiavi in ordine alfabetico
 "Mia" in tel # Restituisce True se presente, False altrimenti
 del tel["Eva"] # Rimuove la coppia chiave-valore con chiave "Eva"
\end{lstlisting}
        \end{minipage}
    \end{alertblock}
\end{frame}

\begin{frame}{METODI PRINCIPALI PER I DIZIONARI}
    \begin{alertblock}{METODI PRINCIPALI}
        \begin{minipage}{0.98\linewidth}
            \justifying
            \begin{itemize}
                \item \textbf{keys()}: Restituisce una lista di tutte le chiavi presenti. 
                \pause
                \item \textbf{values()}: Restituisce una lista di tutti i valori presenti.
                \pause
                \item \textbf{items()}: Restituisce una lista di tutte le coppie (chiave , valore) presenti
            \end{itemize}
        \end{minipage}
    \end{alertblock}
\end{frame}

\end{document}